% !PRINTED 7
\weberpage{7}
%<<CONT
element $a$ belongs to $A$, then every smaller element of $\mathfrak R$
also belongs to $A$. If $A$ is a part of $A'$, we shall call $\alpha$
smaller than $\alpha'$, and thereby the set $\mathfrak S$ has become
ordered.

If we regard the cuts produced by rational fractions as equivalent to
those rational fractions themselves and call them, briefly,
\gesperrt{rational cuts}, then the set $\mathfrak S$ contains the set
$\mathfrak R$, and $\mathfrak S$ is therefore in any case a
\gesperrt{dense set}. But the set $\mathfrak S$ is also
\gesperrt{continuous}. For let us denote the cuts by the large German
letters $\mathfrak A,\mathfrak B,\ldots$, and consider any cut in the
manifold of cuts, $(\mathfrak A,\mathfrak B)$. We can determine an
element of $\mathfrak S$, say $\alpha=(A,B)$, by putting into $A$ every
rational fraction that produces one of the cuts of $\mathfrak A$, and
putting into $B$ all the remaining rational fractions, which therefore
produce the rational cuts in $\mathfrak B$. This cut $\alpha$ in
$\mathfrak R$ produces the cut $(\mathfrak A,\mathfrak B)$ in
$\mathfrak S$. This will be proved when it is shown that every element
$\alpha'$ of $\mathfrak S$ which is smaller than $\alpha$ belongs to
$\mathfrak A$, and every element $\beta'$ of $\mathfrak S$ which is
greater than $\alpha$ belongs to $\mathfrak B$.

Let $\alpha'=(A',B')$ and $\alpha'<\alpha$. Then there are rational
fractions in $A$ which are not contained in $A'$. There is therefore a
rational $\mu$ such that $\alpha'<\mu<\alpha$; this $\mu$ produces a
cut contained in $\mathfrak A$, hence itself belongs to $\mathfrak A$.
Since $\alpha'<\mu$, $\alpha'$ also belongs to $\mathfrak A$. In
exactly the same way one shows that every $\beta'$ greater than
$\alpha$ belongs to $\mathfrak B$, and the continuity of $\mathfrak S$
is thereby proved.

This very abstract way of looking at the matter gives us the assurance
that the assumption of a continuous set contains no contradiction, and
that such sets at least exist in the realm of thought. Geometry, and
also analysis, which has always willingly attached itself to
geometrical intuition, long tacitly assumed the \gesperrt{existence} of
continuous sets, for example in the points of a straight line or of any
other connected line-segment, as a kind of axiom. The
\gesperrt{difference between a dense and a continuous set}, which lies
at the basis of the distinction between commensurable and
incommensurable segments, was also not missed by the ancients.
\setcounter{footnote}{0}\footnote{Euclid, Elements, Book X.}
